\chapter{Topological Spaces}

\begin{center}\itshape
Glue the opposite sides of a square and you get a doughnut; glue them with a twist and you get a surface that no
mirror can see in three dimensions. Remove a point from a line and it falls apart; remove a point from a plane and it
does not. What survives when distances are forgotten, and why is that enough to tell spaces apart?
\end{center}

In 1736 the citizens of K\"onigsberg asked Leonhard Euler whether one could stroll through their town crossing each of its
seven bridges exactly once. Euler answered no, and noticed something odd about his own proof: the lengths of the
bridges, the shape of the river and the size of the islands played no role at all. Only how things were \emph{connected}
mattered. Leibniz had dreamt of such a ``geometry of position''; Euler's little paper was its first theorem, and the
subject later got its name, topology, from Johann Listing (1847).

A topologist, goes the joke, is someone who cannot tell a doughnut from a coffee mug: one can be deformed into the other
without tearing. The joke contains the definition. Topology keeps what survives stretching and bending, and that turns out
to be exactly what continuity needs. Chapter~2 showed that continuity depends only on which sets are open, not on the
distance that defined them. A topology is a declaration of which sets are open, and nothing else; Hausdorff wrote down
the axioms in 1914.

Many natural spaces have open sets but no natural distance: infinite products, spaces glued from squares like the torus and
the projective plane, the spaces of algebraic geometry (Tome~9) and spaces of measures (Chapter~11). This chapter builds them,
with their universal properties. It proves the intermediate value theorem in its natural habitat, connectedness, and meets
a curve that is connected but cannot be walked along. And in the exercises, a topology on the integers proves that there are
infinitely many primes, an argument that Hillel Furstenberg published in 1955 while still an undergraduate.

\section{Topologies}

\begin{definition}
A \emph{topology} on a set $X$ is a family $\mathcal T$ of subsets, the \emph{open sets}, containing $\varnothing$ and $X$ and stable under arbitrary unions
and finite intersections. A set is \emph{closed} if its complement is open. A \emph{neighbourhood} of $x$ is a set containing an open set
containing $x$. The \emph{interior} $\mathring A$ is the largest open set inside $A$, the \emph{closure} $\bar A$ the smallest closed set containing $A$; $A$ is
\emph{dense} if $\bar A=X$. A map $f:X\to Y$ between topological spaces is \emph{continuous} if $f^{-1}(V)$ is open for every open $V$, and a
\emph{homeomorphism} if it is a continuous bijection with continuous inverse.
\end{definition}

So $x\in\bar A$ iff every neighbourhood of $x$ meets $A$ (otherwise an open neighbourhood missing $A$ has a closed complement containing $A$
but not $x$). Compositions of continuous maps are continuous, and a map is continuous iff preimages of closed sets are closed.

\begin{example}\label{ex:topologies}
\begin{enumerate}[label=(\alph*)]
\item The open sets of a metric space (Chapter~2) form the \emph{metric topology}. A topology that comes from a metric is \emph{metrizable}.
\item The \emph{discrete} topology (every set is open) and the \emph{indiscrete} topology ($\{\varnothing,X\}$). Every map out of a discrete space, and every map into an
indiscrete one, is continuous.
\item The \emph{cofinite} topology: the open sets are $\varnothing$ and the complements of finite sets. On an infinite set any two non-empty open sets
meet.
\item The \emph{Sierpi\'nski space} $\{0,1\}$ with open sets $\varnothing$, $\{1\}$, $\{0,1\}$. Continuous maps $X\to\{0,1\}$ are exactly indicators of open sets: the
open sets of $X$ are ``observable'' properties, an idea computer scientists use to model what a program can semi-decide.
\end{enumerate}
\end{example}

If $\mathcal T\subseteq\mathcal T'$, $\mathcal T'$ is \emph{finer}: it has more open sets, so more continuous maps out of it and fewer convergent sequences. Topologies are
usually given by a family of ``basic'' open sets.

\begin{proposition}\label{prop:basis}
Let $\mathcal B$ be a family of subsets of $X$ that covers $X$ and such that, if $x\in B_1\cap B_2$ with $B_1,B_2\in\mathcal B$, then $x\in B_3\subseteq B_1\cap B_2$ for some
$B_3\in\mathcal B$. Then the unions of members of $\mathcal B$ form a topology, the topology \emph{generated} by the \emph{basis} $\mathcal B$. A set $U$ is open for it
iff each $x\in U$ lies in some $B\in\mathcal B$ with $B\subseteq U$, and a map $f$ into $X$ is continuous iff $f^{-1}(B)$ is open for every $B\in\mathcal B$.
\end{proposition}
\begin{proof}
$X$ is a union of members (they cover), $\varnothing$ is the empty union, and unions of unions are unions. If $U=\bigcup B_i$ and $V=\bigcup B'_j$, every
$x\in U\cap V$ is in some $B_i\cap B'_j$, hence in some $B_3\subseteq U\cap V$, so $U\cap V$ is the union of these $B_3$. The characterisation of open sets is a
restatement, and preimages commute with unions.
\end{proof}

The open balls are a basis of a metric topology, the open intervals with rational endpoints a basis of $\R$. More
generally, any family $\mathcal S$ generates a topology, the coarsest containing it: its basis is the finite intersections of members
of $\mathcal S$ (a \emph{subbasis}).

\section{Subspaces, products and quotients}

Each of the three constructions is the coarsest or the finest topology making some maps continuous, and is
characterised by a \emph{universal property}: a statement about all continuous maps into or out of the new space.

\paragraph{Subspaces.} For $A\subseteq X$, the \emph{subspace topology} has open sets $U\cap A$, $U$ open in $X$. It is the coarsest topology making
the inclusion continuous, and a map $f:Z\to A$ is continuous iff it is continuous as a map into $X$ (since $f^{-1}(U\cap A)=f^{-1}(U)$). For a metric
space it is the topology of the restricted metric.

\paragraph{Products.} Let $(X_i)_{i\in I}$ be topological spaces and $X=\prod_iX_i$, the set of families $(x_i)$ with $x_i\in X_i$, with projections $p_i$.

\begin{definition}
The \emph{product topology} on $\prod X_i$ is generated by the \emph{cylinders} $p_i^{-1}(U_i)$, $U_i$ open in $X_i$. Its basic open sets are the finite
intersections $\prod_iU_i$ where $U_i$ is open in $X_i$ and $U_i=X_i$ for all but finitely many $i$.
\end{definition}

\begin{proposition}\label{prop:product}
\begin{enumerate}[label=(\alph*)]
\item A map $f:Z\to\prod X_i$ is continuous iff every $p_i\circ f$ is continuous. A sequence converges iff each coordinate converges.
\item The product topology on $\R^n$ is the topology of any of the norms of Chapter~2.
\item If $(X_k,d_k)_{k\ge1}$ are metric spaces, the product topology on $\prod_kX_k$ is metrizable, by
\[d(x,y)=\sum_{k\ge1}2^{-k}\min\bigl(1,d_k(x_k,y_k)\bigr).\]
\end{enumerate}
\end{proposition}
\begin{proof}
(a) If $f$ is continuous, so is $p_i\circ f$. Conversely the preimage of a cylinder is $f^{-1}(p_i^{-1}(U_i))=(p_i\circ f)^{-1}(U_i)$, open, and cylinders generate.
A sequence $x^{(n)}$ converges to $x$ iff it eventually enters every basic open set around $x$; such a set constrains finitely many
coordinates, so this holds iff each coordinate converges.
(b) An open box $\prod\,]a_i,b_i[$ is open for $\|\cdot\|_\infty$ and contains a $\|\cdot\|_\infty$-ball around each of its points, and a $\|\cdot\|_\infty$-ball is an open box; so
the two topologies have the same open sets, and the norms of Chapter~2 are Lipschitz equivalent.
(c) $d$ is a metric (each term is, and $\min(1,\cdot)$ keeps the triangle inequality). Let $B$ be a $d$-ball around $x$ of radius $r$: choose $K$ with
$2^{-K}<\frac r2$; the basic set $\prod_{k\le K}B_{d_k}(x_k,\frac r2)\times\prod_{k>K}X_k$ lies in $B$, since $d(x,y)<\sum_{k\le K}2^{-k}\frac r2+2^{-K}<r$. Conversely a basic open set
around $x$ constrains finitely many coordinates, each to contain a ball $B_{d_k}(x_k,\varepsilon)$ with $\varepsilon<1$, and the $d$-ball of radius $2^{-K}\varepsilon$, with
$K$ the largest constrained index, lies inside it, because $d(x,y)<2^{-K}\varepsilon$ forces $d_k(x_k,y_k)<\varepsilon$ for $k\le K$.
\end{proof}

The product $\{0,1\}^\N$ (the \emph{Cantor space}) and $[0,1]^\N$ (the \emph{Hilbert cube}) are therefore metrizable. An uncountable product of
spaces with two or more points is not (Exercise~\ref{ex:cocountable} gives the flavour: sequences do not suffice there).

\paragraph{Quotients.} Let $\sim$ be an equivalence relation on $X$ and $\pi:X\to X/{\sim}$ the quotient map.

\begin{definition}
The \emph{quotient topology} on $X/{\sim}$ declares $V$ open iff $\pi^{-1}(V)$ is open in $X$. It is the finest topology making $\pi$ continuous.
\end{definition}

\begin{proposition}[universal property of quotients]\label{prop:quotient}
A map $g:X/{\sim}\to Z$ is continuous iff $g\circ\pi$ is continuous. So a continuous $f:X\to Z$ that is constant on classes factors as $f=\bar f\circ\pi$
with $\bar f$ continuous.
\end{proposition}
\begin{proof}
$(g\circ\pi)^{-1}(W)=\pi^{-1}(g^{-1}(W))$, which is open iff $g^{-1}(W)$ is open, by definition of the quotient topology.
\end{proof}

\begin{figure}[ht]
\centering
\fig{ch03-gluing}
\caption{Surfaces from a square $[0,1]^2$ by gluing edges as the arrows indicate: the torus (both pairs of sides glued with the same
orientation), the Klein bottle (one pair reversed) and the projective plane (both pairs reversed). The Klein bottle and the projective plane
cannot be embedded in $\R^3$.}
\label{fig:gluing}
\end{figure}

\begin{example}\label{ex:quotients}
\begin{enumerate}[label=(\alph*)]
\item The circle. $\R/\Z$ (identify $x$ and $y$ when $x-y\in\Z$) maps to $S^1=\{z\in\C:|z|=1\}$ by $\bar f([x])=e^{2i\pi x}$, a continuous bijection by
Proposition~\ref{prop:quotient} and Theorem~\ref{thm:exp}(e). Its inverse is continuous too, because $\pi$ is \emph{open} (the image of an open $U$ has
preimage $\bigcup_{n\in\Z}(U+n)$, open) and $f$ maps small open intervals onto open arcs. So $\R/\Z\cong S^1$.
\item The torus. $\R^2/\Z^2\cong S^1\times S^1$ by the same argument coordinatewise, and the square with opposite sides glued is the same
space (Chapter~4 shows that a continuous bijection from a compact space to a Hausdorff one is a homeomorphism, which makes such
identifications automatic). Pac-Man lives on this torus.
\item A bad quotient. On $\R/\Q$ (identify $x$ and $y$ when $x-y\in\Q$), a non-empty open $V$ has an open, non-empty preimage stable under rational
translations, hence dense; its complement is closed, stable under rational translations, with empty interior. One checks
that such a non-empty open set must be all of $\R$ (Exercise~\ref{ex:RmodQ}), so $\R/\Q$ is indiscrete: the quotient remembers nothing. This
is the topological face of Vitali's non-measurable set (Chapter~6).
\end{enumerate}
\end{example}

\section{Connectedness}

In 1817 Bernard Bolzano, a priest in Prague, published a short paper with a long title: \emph{Purely analytic proof of the
theorem that between any two values which give results of opposite sign, there lies at least one real root of the
equation.} Everyone ``knew'' that a continuous curve going from below the axis to above it must cross it; Bolzano insisted
that a picture is not a proof, and that the reason lies in the real numbers themselves. He was right, and the reason has a
name: the line is \emph{connected}, it cannot be split into two open pieces.

\begin{definition}
A space $X$ is \emph{connected} if it is not the union of two disjoint non-empty open sets; equivalently, if the only subsets that are both
open and closed are $\varnothing$ and $X$; equivalently, if every continuous map $X\to\{0,1\}$ (discrete) is constant. A subset is connected if it is
connected as a subspace.
\end{definition}

The equivalences: a set $U$ that is open and closed gives the partition $X=U\cup(X\setminus U)$ into open sets, and the indicator of $U$ is a
continuous map to $\{0,1\}$; conversely the preimages of $0$ and $1$ under a continuous map are open, disjoint, and cover $X$.

\begin{theorem}\label{thm:intervals}
A subset of $\R$ is connected iff it is an interval (a set $I$ such that $a<c<b$ with $a,b\in I$ implies $c\in I$).
\end{theorem}
\begin{proof}
If $I$ is not an interval, there are $a<c<b$ with $a,b\in I$ and $c\notin I$; then $I\cap\,]-\infty,c[$ and $I\cap\,]c,+\infty[$ separate $I$. Conversely let $I$ be an
interval and $f:I\to\{0,1\}$ continuous; suppose $f(a)=0$ and $f(b)=1$ with $a<b$ (the case $a>b$ is symmetric). Let $s=\sup\{x\in[a,b]:f=0\text{ on }[a,x]\}$; then
$s\in[a,b]\subseteq I$. By continuity $f$ is constant on a neighbourhood $]s-\eta,s+\eta[\,\cap\,[a,b]$ of $s$ in $[a,b]$. If $f(s)=0$, then $f=0$ on $[a,s]$ (points left of $s$ are
in some $[a,x]$ with $f=0$, and $f(s)=0$), hence on $[a,\min(s+\eta/2,b)]$, so $s=b$ by maximality, contradicting $f(b)=1$. If $f(s)=1$, then $s>a$
and $f=1$ on $]s-\eta,s]$, but $f=0$ on $[a,x]$ for some $x>s-\eta$: a contradiction.
\end{proof}

\begin{proposition}\label{prop:connected}
\begin{enumerate}[label=(\alph*)]
\item A continuous image of a connected space is connected.
\item A union of connected subsets that all contain a common point is connected; the closure of a connected subset is connected.
\item A product of connected spaces is connected.
\end{enumerate}
\end{proposition}
\begin{proof}
(a) If $f:X\to Y$ is onto and $g:Y\to\{0,1\}$ continuous, $g\circ f$ is constant, so $g$ is. (b) A continuous $g:\bigcup A_i\to\{0,1\}$ is constant on each $A_i$,
with the value it takes at the common point. If $A$ is connected and $g:\bar A\to\{0,1\}$ is continuous, $g$ is constant on $A$, equal to $c$; then
$g^{-1}(c)$ is closed in $\bar A$ and contains $A$, so it is $\bar A$. (c) For two factors: fix $(a,b)\in X\times Y$. For each $x$, the set
$(\{x\}\times Y)\cup(X\times\{b\})$ is a union of two connected sets (homeomorphic to $Y$ and $X$) meeting at $(x,b)$, hence connected; these sets all
contain $(a,b)$ and cover $X\times Y$; apply (b). By induction, finite products are connected. For an arbitrary product, fix $a\in\prod X_i$; the
points that differ from $a$ in finitely many coordinates form a union of finite products (each homeomorphic to a finite product of
factors, times one point) all containing $a$, so a connected set $D$. It is dense: a basic open set constrains finitely many coordinates,
and contains a point of $D$. By (b), $\prod X_i=\bar D$ is connected.
\end{proof}

\begin{corollary}[intermediate value theorem]\label{cor:ivt}
If $X$ is connected and $f:X\to\R$ is continuous, then $f(X)$ is an interval. In particular a continuous function on an interval takes
every value between two of its values.
\end{corollary}
\begin{proof}
$f(X)$ is connected by Proposition~\ref{prop:connected}(a), hence an interval by Theorem~\ref{thm:intervals}.
\end{proof}

This is the theorem Tome~0 admitted, and Lemma~\ref{lem:monoivt} of Chapter~2 was a special case. The \emph{connected components} of $X$ are the
maximal connected subsets: by Proposition~\ref{prop:connected}(b), the union of all connected sets containing $x$ is connected, and these
unions partition $X$ into closed sets. The components of $\Q$ are its points, and so are those of the Cantor set: such spaces are
\emph{totally disconnected}.

\begin{definition}
$X$ is \emph{path-connected} if any two points are joined by a \emph{path}, a continuous $\gamma:[0,1]\to X$ with $\gamma(0)=x$, $\gamma(1)=y$.
\end{definition}

A path-connected space is connected: all the path images $\gamma([0,1])$ from a fixed point are connected (Theorem~\ref{thm:intervals} and
Proposition~\ref{prop:connected}(a)) and contain that point. Convex subsets of $\R^n$ are path-connected (segments), and so are spheres
$S^{n-1}$ for $n\ge2$ (go along great circles).

\begin{proposition}\label{prop:openconnected}
\begin{enumerate}[label=(\alph*)]
\item An open connected subset $U$ of a normed space is path-connected, even by polygonal paths.
\item The \emph{topologist's sine curve} $T=\{(x,\sin\frac1x):0<x\le1\}\cup(\{0\}\times[-1,1])$ is connected but not path-connected.
\end{enumerate}
\end{proposition}
\begin{proof}
(a) Fix $a\in U$ and let $P$ be the set of points of $U$ joined to $a$ by a polygonal path in $U$. If $x\in U$, a ball $B(x,r)\subseteq U$ is convex, so either
all of it is in $P$ (if $x\in P$, or any point of the ball is) or none of it is. So $P$ and $U\setminus P$ are open; $P\ni a$ is non-empty, hence $P=U$.
(b) The graph part $G$ is connected (a continuous image of $]0,1]$) and $T=\bar G$, since $\sin\frac1x$ takes every value of $[-1,1]$ on every interval
$]0,\varepsilon[$; so $T$ is connected by Proposition~\ref{prop:connected}(b). Suppose $\gamma=(\gamma_1,\gamma_2)$ is a path in $T$ from $(0,0)$ to $(1,\sin1)$. Let
$t_0=\sup\{t:\gamma_1(t)=0\}$; by continuity $\gamma_1(t_0)=0$, and $\gamma_1(t)>0$ for $t>t_0$. By continuity of $\gamma_2$ at $t_0$ there is $\delta>0$ with
$|\gamma_2(t)-\gamma_2(t_0)|<\frac12$ on $[t_0,t_0+\delta]$. But $\gamma_1([t_0,t_0+\delta])$ is an interval $[0,c]$ with $c>0$ (Corollary~\ref{cor:ivt}), and on it $\sin\frac1x$ takes
both values $1$ and $-1$, at points $x=\gamma_1(t)$ with $\gamma_2(t)=\sin\frac1{\gamma_1(t)}$. Both cannot be within $\frac12$ of $\gamma_2(t_0)$.
\end{proof}

\begin{figure}[ht]
\centering
\fig{ch03-sine}
\caption{The topologist's sine curve: the graph of $\sin\frac1x$ on $]0,1]$ (blue) accumulates on the segment $\{0\}\times[-1,1]$ (orange). The whole set is
connected, but no path can cross from the segment to the graph: near $x=0$ the graph oscillates infinitely often.}
\label{fig:sine}
\end{figure}

\begin{think}
Is $\R$ homeomorphic to $\R^2$? Both are path-connected. What happens to each when one point is removed?
\end{think}

Removing a point disconnects $\R$ but not $\R^2$ (a punctured plane is path-connected), and a homeomorphism $\R\to\R^2$ would restrict to a
homeomorphism between the punctured spaces. So $\R\not\cong\R^2$ (Exercise~\ref{ex:notR2}). That $\R^2\not\cong\R^3$ is harder: it needs the fundamental
group or homology of Tome~8 (Brouwer's invariance of domain).

\section{Separation and sequences}

Topology is generous: it allows spaces so coarse that two different points cannot be told apart by any open set, and
sequences that converge to several limits at once. Such spaces do occur (in algebraic geometry, for instance), but analysis
needs a minimum of politeness: distinct points should have disjoint neighbourhoods. Hausdorff put this axiom into his very
definition of a topological space in 1914; it now carries his name.

\begin{definition}
A space is \emph{Hausdorff} if two distinct points have disjoint neighbourhoods; \emph{normal} if moreover points are closed and two disjoint
closed sets have disjoint open neighbourhoods.
\end{definition}

In a Hausdorff space a sequence has at most one limit, and $\{x\}$ is closed. Metric spaces are Hausdorff (balls of radius
$\frac{d(x,y)}2$). Subspaces and products of Hausdorff spaces are Hausdorff (basic open sets separate a differing coordinate). Quotients need not
be: $\R/\Q$ is not, and neither is the \emph{line with two origins}, $\R\times\{0,1\}$ with $(x,0)\sim(x,1)$ for $x\ne0$, where the two origins cannot be
separated.

\begin{proposition}[Urysohn functions in metric spaces]\label{prop:urysohn}
Let $A$, $B$ be disjoint closed subsets of a metric space $X$. Then
\[
f(x)=\frac{d(x,A)}{d(x,A)+d(x,B)}
\]
is continuous $X\to[0,1]$, equal to $0$ exactly on $A$ and to $1$ exactly on $B$. Hence metric spaces are normal.
\end{proposition}
\begin{proof}
$d(x,A)=0$ iff $x\in\bar A=A$ (Proposition~\ref{prop:closure}), and likewise for $B$; as $A\cap B=\varnothing$, the denominator never vanishes. The distance
functions are $1$-Lipschitz, so $f$ is continuous. The open sets $\{f<\frac12\}$ and $\{f>\frac12\}$ separate $A$ and $B$.
\end{proof}

\begin{remark}
Urysohn's lemma (1925) says that in any normal space there is such a continuous $f$ (with $A\subseteq f^{-1}(0)$ and $B\subseteq f^{-1}(1)$), although there is no
formula; its proof builds the level sets $\{f<r\}$ for dyadic $r$ one by one (J.~Munkres, \emph{Topology}, \S33). Compact Hausdorff spaces are
normal (Chapter~4). The series only needs the metric case.
\end{remark}

In a metric space, sequences decide everything (Proposition~\ref{prop:closure}). The reason is that each point has a countable family of
neighbourhoods $B(x,\frac1n)$ that are ``small enough''.

\begin{definition}
$X$ is \emph{first countable} if each point has a countable family of neighbourhoods $V_1,V_2,\ldots$ such that every neighbourhood contains
some $V_n$; \emph{second countable} if the topology has a countable basis; \emph{separable} if it has a countable dense subset.
\end{definition}

\begin{proposition}\label{prop:sequences}
\begin{enumerate}[label=(\alph*)]
\item In a first countable space, $x\in\bar A$ iff some sequence of $A$ converges to $x$, and $f:X\to Y$ is continuous iff it preserves limits of
sequences.
\item A metric space is second countable iff it is separable. $\R^n$, $C([0,1])$ (sup norm, by Weierstrass, Chapter~4) and $\prod_k X_k$ for separable metric
$X_k$ are separable.
\end{enumerate}
\end{proposition}
\begin{proof}
(a) Replacing $V_n$ by $V_1\cap\cdots\cap V_n$ we may assume them decreasing. If $x\in\bar A$, pick $a_n\in A\cap V_n$; every neighbourhood contains some $V_N$,
hence all $a_n$ for $n\ge N$: $a_n\to x$. Conversely a limit of points of $A$ has every neighbourhood meeting $A$. If $f$ preserves limits and $F$ is
closed in $Y$, let $x\in\overline{f^{-1}(F)}$ and $a_n\to x$ with $f(a_n)\in F$; then $f(x)=\lim f(a_n)\in F$. So $f^{-1}(F)$ is closed. The converse holds in any space.
(b) If $\{U_n\}$ is a countable basis, one point from each non-empty $U_n$ is dense. If $D=\{d_k\}$ is dense, the balls $B(d_k,\frac1m)$ form a basis:
for $x\in U$ open, take $B(x,\frac2m)\subseteq U$ and $d_k\in B(x,\frac1m)$; then $x\in B(d_k,\frac1m)\subseteq B(x,\frac2m)\subseteq U$. For $\R^n$ take $\Q^n$; for $C([0,1])$ the
polynomials with rational coefficients, dense by Weierstrass's theorem (Chapter~4); for a countable product, the points with
coordinates in the countable dense sets $D_k$ and equal to a fixed $a_k$ for all large $k$ (a countable set by Proposition~\ref{prop:countable},
dense by Proposition~\ref{prop:product}).
\end{proof}

When neighbourhoods are not countable, sequences may fail to detect closures (Exercise~\ref{ex:cocountable}). The fix is to use \emph{nets}
(families indexed by directed sets) or \emph{filters}; this tome avoids them, and proves Tychonoff's theorem directly in Chapter~4.

\begin{beauty}
Topology keeps exactly the information needed for continuity and forgets the rest, and this is why it applies so
widely: to number theory (Furstenberg's primes, Exercise~\ref{ex:furstenberg}, and the profinite integers), to algebraic geometry,
where the Zariski topology has very few open sets but still tells irreducible varieties apart, to logic and computer
science, where open sets are semi-decidable properties, and to geometry, where gluing squares produces every closed
surface. The intermediate value theorem, admitted in Tome~0, turns out to be a statement about connectedness of
intervals, which is in turn a reformulation of the least upper bound property of $\R$.
\end{beauty}

\section{Exercises}

\begin{exercise}[$\star$]\label{ex:cofinite}
Let $X$ be an infinite set with the cofinite topology. Show that $X$ is connected, that every map $X\to X$ with finite fibres is continuous, and that a
sequence of distinct points converges to every point of $X$. Is $X$ Hausdorff?
\end{exercise}

\begin{exercise}[$\star$]\label{ex:RmodQ}
Show that the quotient $\R/\Q$ of Example~\ref{ex:quotients}(c) is indiscrete: a non-empty open subset of $\R$ that is stable under all rational
translations is $\R$. (Use that it is open and that $\Q$ is dense.)
\end{exercise}

\begin{exercise}[$\star\star$]\label{ex:notR2}
Show that $\R$ and $\R^2$ are not homeomorphic, and that $[0,1]$ and $[0,1]^2$ are not homeomorphic. Show also that the circle $S^1$ is not homeomorphic to
an interval.
\end{exercise}

\begin{exercise}[$\star\star$]\label{ex:monotone}
Let $I$ be an interval and $f:I\to\R$ continuous and injective. Show that $f$ is strictly monotone, that $J=f(I)$ is an interval, and that $f^{-1}:J\to I$ is
continuous.
\end{exercise}

\begin{exercise}[$\star\star$]\label{ex:furstenberg}
\emph{Furstenberg revisited.} For $n\in\Z\setminus\{0\}$ let $v(n)$ be the largest $k$ such that $k!$ divides $n$, and $v(0)=+\infty$. (a) Show that $d(m,n)=2^{-v(m-n)}$
is an ultrametric on $\Z$ whose topology is the one of Tome~1, Exercise~2.8: the open sets are the unions of arithmetic progressions $a+b\Z$,
$b\ge1$. (b) Show that every progression $a+b\Z$ is open and closed, so that $\Z$ is totally disconnected. (c) Deduce Furstenberg's proof: if there
were finitely many primes, $\{-1,1\}$ would be open.
\end{exercise}

\begin{exercise}[$\star\star$]\label{ex:finite}
Show that topologies on a finite set $X$ correspond bijectively to preorders on $X$ (reflexive transitive relations), by $x\preceq y$ iff $x\in\overline{\{y\}}$.
Which preorders give $T_0$ spaces (two distinct points are distinguished by an open set)? How many topologies are there on a set with
three elements?
\end{exercise}

\begin{exercise}[$\star\star$]\label{ex:cocountable}
On $\R$, declare open $\varnothing$ and the complements of countable sets (the \emph{cocountable} topology). Show that a sequence converges only if it is
eventually constant, that $0$ is in the closure of $A=\R\setminus\{0\}$, but that no sequence of $A$ converges to $0$. Deduce that this space is not first
countable, and that the identity from it to $\R$ with the discrete topology preserves limits of sequences without being continuous.
\end{exercise}

\begin{exercise}[$\star\star\star$]\label{ex:tietze}
\emph{Tietze's extension theorem for metric spaces.} Let $A$ be closed in a metric space $X$ and $f:A\to[-1,1]$ continuous. (a) Show that there is a
continuous $g:X\to[-\frac13,\frac13]$ with $|f-g|\le\frac23$ on $A$. (b) Iterate to build $g_1,g_2,\ldots$ with $|f-\sum_{k\le n}g_k|\le(\frac23)^n$ on $A$ and
$|g_k|\le\frac13(\frac23)^{k-1}$, and conclude that $f$ extends to a continuous $F:X\to[-1,1]$.
\end{exercise}

\section{Hints}

\begin{hint}{ex:cofinite}
Two non-empty open sets always meet. A neighbourhood of $x$ misses only finitely many points.
\end{hint}
\begin{hint}{ex:RmodQ}
If $U$ is open, stable under $\Q$ and contains $x$, it contains an interval around $x$, and its translates by rationals.
\end{hint}
\begin{hint}{ex:notR2}
Remove one point (for the circle, remove two). A homeomorphism restricts to a homeomorphism of the complements.
\end{hint}
\begin{hint}{ex:monotone}
If $a<b<c$ in $I$ and $f(b)$ is not between $f(a)$ and $f(c)$, the intermediate value theorem gives two points with the same image. For the
inverse, a strictly monotone map onto an interval is open.
\end{hint}
\begin{hint}{ex:furstenberg}
$v(m+n)\ge\min(v(m),v(n))$. The ball of radius $2^{-k}$ around $a$ is $a+(k+1)!\,\Z$, and $a+b\Z$ is a union of such balls when $b\mid k!$.
\end{hint}
\begin{hint}{ex:finite}
Recover the open sets as the sets closed upwards for $\preceq$. Count preorders on three points by their shapes.
\end{hint}
\begin{hint}{ex:cocountable}
The set of terms different from the limit is countable, hence closed.
\end{hint}
\begin{hint}{ex:tietze}
(a) Apply Proposition~\ref{prop:urysohn} to $\{f\le-\frac13\}$ and $\{f\ge\frac13\}$, closed in $X$ because $A$ is closed.
\end{hint}

\section{Solutions}

\begin{solution}{ex:cofinite}
\idea Open sets are huge: each one misses only finitely many points. So they all meet, and neighbourhoods are hard to avoid.
\writeup Two non-empty open sets $X\setminus F_1$, $X\setminus F_2$ meet, since $X$ is infinite and $F_1\cup F_2$ finite; so $X$ is not a disjoint union of two non-empty open sets, and is
connected, and not Hausdorff. If $g:X\to X$ has finite fibres, the preimage of a finite set is finite, so preimages of closed sets are closed: $g$ is
continuous. If $(x_n)$ are distinct and $U=X\setminus F$ is a neighbourhood of $x$, at most $|F|$ terms lie in $F$, so $x_n\in U$ for all large $n$: $x_n\to x$, for every $x$.
\end{solution}

\begin{solution}{ex:RmodQ}
\idea Openness gives an interval; stability under $\Q$ spreads that interval over the whole line, because its rational translates
cover $\R$.
\writeup Let $U\subseteq\R$ be open, non-empty, with $U+q=U$ for all $q\in\Q$, and $x\in U$, $]x-\varepsilon,x+\varepsilon[\,\subseteq U$. For $y\in\R$, pick $q\in\Q$ with $|y-x-q|<\varepsilon$
(density). Then $y-q\in\,]x-\varepsilon,x+\varepsilon[\,\subseteq U$, so $y\in U+q=U$. Hence $U=\R$. The open sets of $\R/\Q$ are the $V$ whose preimage is open, and that preimage
is $\Q$-stable; so $V=\varnothing$ or $V$ is everything.
\end{solution}

\begin{solution}{ex:notR2}
\idea Connectedness is preserved by homeomorphisms, and so is the effect of removing points; count what happens.
\writeup If $h:\R\to\R^2$ were a homeomorphism, it would restrict to a homeomorphism $\R\setminus\{0\}\to\R^2\setminus\{h(0)\}$. The first is not connected (Theorem~\ref{thm:intervals}),
the second is path-connected (join two points by a polygonal path with one corner, avoiding the removed point), hence connected:
impossible. The same argument works for $[0,1]$ and $[0,1]^2$, removing the point $\frac12$ and its image (a square minus a point is path-connected,
by the same polygonal paths). For the circle: removing any one point of $S^1$ leaves a connected set (an arc, the image of an interval),
while removing an interior point of an interval disconnects it. If $h:S^1\to I$ were a homeomorphism, the interval $I$ has at least two points,
hence an interior point $c$, and $S^1\setminus\{h^{-1}(c)\}$ would be homeomorphic to the disconnected $I\setminus\{c\}$.
\end{solution}

\begin{solution}{ex:monotone}
\idea Injective and continuous on an interval forbids ``turning back'', since turning back would hit some value twice by the intermediate
value theorem.
\writeup Let $a<b<c$ in $I$. If $f(b)$ were not strictly between $f(a)$ and $f(c)$, say $f(b)>\max(f(a),f(c))$, then any value $y$ with $\max(f(a),f(c))<y<f(b)$ is taken
on $]a,b[$ and on $]b,c[$ (Corollary~\ref{cor:ivt}), contradicting injectivity; similarly if $f(b)<\min$. So $f$ is strictly monotone on every triple. Then it is
monotone: if $f(a)<f(c)$ for some $a<c$ and $f(x)>f(y)$ for some $x<y$, applying the triple property to suitable triples among $a,c,x,y$
gives a contradiction (fix $a<c$, and slide $x,y$ by triples to $a,c$). So $f$ is, say, increasing. $J=f(I)$ is an interval by Corollary~\ref{cor:ivt}. The inverse
is increasing; for continuity at $y_0=f(x_0)$ and $\varepsilon>0$ with $x_0\pm\varepsilon\in I$ (otherwise use one side), $f(x_0-\varepsilon)<y_0<f(x_0+\varepsilon)$, and for $y$ between these two
values, $f^{-1}(y)$ is between $x_0-\varepsilon$ and $x_0+\varepsilon$.
\end{solution}

\begin{solution}{ex:furstenberg}
\idea The factorials are divisible by every integer from some point on, so balls for $d$ are progressions with factorial differences, and
every progression is a union of them. In this topology non-empty open sets are infinite, while finite unions of progressions are closed.
\writeup (a) If $k!$ divides $m$ and $n$, it divides $m+n$: $v(m+n)\ge\min(v(m),v(n))$, so $d$ is an ultrametric ($d(m,n)=0$ iff $m=n$, since a non-zero integer has finite
$v$). The open ball $B(a,2^{-k})=\{n:v(n-a)>k\}$; since $(k+1)!$ divides $j!$ for $j\ge k+1$, this is $\{n:(k+1)!\mid n-a\}=a+(k+1)!\Z$. If $n\in a+b\Z$ and $b\mid k!$, then
$n+k!\Z\subseteq a+b\Z$: progressions are open for $d$; and $d$-balls are progressions. So the two topologies coincide. (b) $a+b\Z$ is open, and its
complement is the union of the progressions $a+j+b\Z$, $1\le j<b$, also open. Two distinct integers are separated by a progression and its
complement, both open, so components are points. (c) Every integer other than $\pm1$ has a prime factor, so $\Z\setminus\{-1,1\}=\bigcup_pp\Z$. If there
were finitely many primes, this would be a finite union of closed sets, hence closed, and $\{-1,1\}$ would be open. But a non-empty open set
contains a progression, which is infinite.
\end{solution}

\begin{solution}{ex:finite}
\idea On a finite set every intersection of open sets is finite, so each point has a smallest open neighbourhood $U_x$; the topology is determined
by these, and ``$y\in U_x$'' is a preorder.
\writeup For a topology, let $U_x$ be the intersection of the (finitely many) open sets containing $x$, an open set. Then $x\in\overline{\{y\}}$ iff every open set
containing $x$ contains $y$, iff $y\in U_x$. This relation is reflexive, and transitive ($y\in U_x$ and $z\in U_y$ give $U_y\subseteq U_x$, so $z\in U_x$). A set $V$ is open iff
it contains $U_x$ for each $x\in V$, iff it is closed upwards ($x\in V$ and $x\preceq y$ imply $y\in V$). Conversely, for a preorder the upward-closed sets are stable
under unions and intersections, form a topology, and the closure of $\{y\}$ is $\{x:x\preceq y\}$, the complement of the largest upward-closed set
missing $y$; so the two constructions are inverse. The space is $T_0$ iff $\preceq$ is antisymmetric (a partial order): $x\preceq y\preceq x$ means $x$, $y$ lie in the
same open sets. On three points there are $29$ topologies (up to relabelling, $9$): count preorders by the partition into classes
$x\preceq y\preceq x$ and a partial order on the classes; for instance $19$ partial orders on three labelled points, $9$ preorders with one class of two
points and one single point, and $1$ with all three equivalent.
\end{solution}

\begin{solution}{ex:cocountable}
\idea Countable sets are closed, so the tail of a sequence that avoids its limit is a closed set that a neighbourhood can dodge. The
closure of a large set, on the contrary, is everything.
\writeup If $x_n\to x$, the set $S$ of terms different from $x$ is countable, so $U=\R\setminus S$ is an open neighbourhood of $x$, and $x_n\in U$ means $x_n=x$ for all large $n$.
Every non-empty open set is cocountable, hence meets $A$ (uncountable): $0\in\bar A$. But a sequence of $A$ converging to $0$ would be eventually
equal to $0\notin A$. By Proposition~\ref{prop:sequences}(a), the space is not first countable. The identity to the discrete topology preserves limits
(only eventually constant sequences converge, and they converge in any topology) but is not continuous: $\{0\}$ is open in the
discrete topology, not in the cocountable one.
\end{solution}

\begin{solution}{ex:tietze}
\idea Urysohn's function approximates $f$ to within two thirds of its size; the error is again a continuous function on $A$, smaller by a
factor $\frac23$, and the corrections form a uniformly convergent series.
\writeup (a) Let $B=\{x\in A:f(x)\le-\frac13\}$ and $C=\{x\in A:f(x)\ge\frac13\}$, closed in $A$ hence in $X$, and disjoint. If both are non-empty, let $u$ be the function of
Proposition~\ref{prop:urysohn} ($0$ on $B$, $1$ on $C$) and $g=\frac23u-\frac13$, with values in $[-\frac13,\frac13]$; if $B$ (or $C$) is empty, take $g=\frac13$ (or $-\frac13$), and $g=0$ if
both are. On $B$, $g=-\frac13$ and $f\in[-1,-\frac13]$; on $C$, $g=\frac13$ and $f\in[\frac13,1]$; elsewhere on $A$, $f$ and $g$ are in $[-\frac13,\frac13]$. In all cases
$|f-g|\le\frac23$. (b) Apply (a), scaled by $(\frac23)^n$, to $f-\sum_{k\le n}g_k$, which takes values in $[-(\frac23)^n,(\frac23)^n]$ on $A$: this gives $g_{n+1}$ with
$|g_{n+1}|\le\frac13(\frac23)^n$ and $|f-\sum_{k\le n+1}g_k|\le(\frac23)^{n+1}$ on $A$. The series $F=\sum_kg_k$ converges uniformly on $X$ (geometric bound), so $F$ is continuous
(Proposition~\ref{prop:complete}(c)), $|F|\le\frac13\sum_k(\frac23)^{k-1}=1$, and $F=f$ on $A$.
\end{solution}

\section*{Sources and further reading}

\begin{itemize}
\item J.~Munkres, \emph{Topology}, 2nd ed., Prentice Hall, 2000, Chapters~2--4. The standard first course, with products, quotients,
connectedness and the countability and separation axioms, and the proof of Urysohn's lemma.
\item F.~Hausdorff, \emph{Grundz\"uge der Mengenlehre}, Veit, Leipzig, 1914. Topological spaces defined by neighbourhoods, and the
separation axiom that carries his name.
\item L.~A.~Steen and J.~A.~Seebach, \emph{Counterexamples in Topology}, 2nd ed., Springer, 1978 (Dover, 1995). The cofinite,
cocountable and Sierpi\'nski spaces, the topologist's sine curve and 139 more.
\item H.~Furstenberg, \emph{On the infinitude of primes}, American Mathematical Monthly 62 (1955), 353.
\item S.~Vickers, \emph{Topology via Logic}, Cambridge University Press, 1989. Open sets as observable properties, for computer
scientists.
\item J.~Stillwell, \emph{Classical Topology and Combinatorial Group Theory}, 2nd ed., Springer, 1993. Surfaces as glued polygons.
\end{itemize}
